\chapter{Cervical Cerclage}

\section*{Introduction \& Clinical Indications}
Cervical cerclage is the placement of a suture around the cervix to help delay painless cervical dilation and reduce the risk of pregnancy loss or spontaneous preterm birth in carefully selected pregnancies. It is not a treatment for established preterm labor, intra-amniotic infection, or preterm prelabor rupture of membranes (PPROM). The decision is individualized by an obstetrician, often with maternal--fetal medicine input.

Cerclage may be considered in three settings:
\begin{itemize}
  \item \textbf{History-indicated cerclage}: a prior second-trimester loss or very early spontaneous preterm birth associated with painless cervical dilation, or a prior cerclage for that indication. It is usually placed in the late first or early second trimester.
  \item \textbf{Ultrasound-indicated cerclage}: cervical shortening on standardized transvaginal ultrasound in a patient with relevant prior spontaneous preterm birth risk. A short cervix alone in an asymptomatic singleton pregnancy without that history is generally managed with vaginal progesterone rather than cerclage.
  \item \textbf{Exam-indicated (rescue) cerclage}: painless cervical dilation in the second trimester after assessment has excluded labor, PPROM, and intra-amniotic infection. Evidence is more limited, so counseling is particularly important.
\end{itemize}

Transabdominal cerclage is an alternative for selected patients in whom a transvaginal cerclage is not technically feasible or who had a very early delivery despite a prior transvaginal cerclage. It requires cesarean delivery while the stitch remains in place.

\section*{Relevant Surgical Anatomy}
The cervix is the lower portion of the uterus and surrounds the endocervical canal. A transvaginal stitch is placed high around the cervix near the cervicovaginal junction, while avoiding the bladder anteriorly and the parametrial tissues and uterine vessels laterally. The McDonald technique is a purse-string suture placed without dissecting a bladder flap; the Shirodkar technique is placed higher and may require limited mucosal dissection. The cervix is supplied mainly by branches of the uterine and vaginal arteries, so bleeding control is important.

\section*{Preoperative Workup \& Preparation}
Before placement, the indication, gestational age, fetal viability and anatomy, cervical length or dilation, and alternative management options are reviewed. Evaluation is directed at excluding uterine contractions, PPROM, significant bleeding, placental complications, and clinical or laboratory evidence of intra-amniotic infection. Testing may include a speculum examination, fetal assessment, transvaginal ultrasound when appropriate, and targeted laboratory or infection testing; routine broad laboratory panels are not required for every patient.

Informed consent should cover possible failure to prevent preterm birth, PPROM, infection, bleeding, cervical laceration, anesthesia complications, and the possibility of emergency stitch removal or delivery. Regional or general anesthesia may be used. Food and fluid restrictions follow the anesthesia service's instructions. Routine prophylactic antibiotics and routine tocolysis are not universally required and should be based on the clinical situation and local protocol.

Cerclage is generally avoided when there is active labor, significant unexplained bleeding, PPROM, suspected or confirmed intra-amniotic infection, or a fetal condition for which continuing the pregnancy is not appropriate. These are clinical contraindications rather than absolute rules; urgent specialist assessment is required.

\section*{Surgical Technique \& Procedure}
For a transvaginal cerclage, the patient is placed in lithotomy position and receives regional or general anesthesia. A speculum exposes the cervix, which is cleansed. In the McDonald technique, a strong suture is passed circumferentially through the cervical stroma in several bites, avoiding the cervical canal, and tied to provide support without excessive tissue strangulation. The knot is left accessible for later removal. A Shirodkar stitch or transabdominal approach may be selected when a higher stitch is needed or vaginal placement is not feasible. Fetal heart activity and maternal recovery are assessed after the procedure.

\section*{Complications \& Risks}
Potential complications include bleeding, infection of the cervix or fetal membranes, PPROM, uterine contractions or preterm labor, cervical laceration, stitch displacement or failure, and anesthesia-related complications. If labor or PPROM occurs with a transvaginal stitch in place, the stitch usually needs prompt assessment and often removal to reduce the risk of cervical injury. Rarely, hemorrhage or injury to adjacent organs occurs.

\section*{Postoperative Care \& Recovery}
Most uncomplicated transvaginal cerclages are managed as day surgery or with a short observation period. Mild cramping, spotting, or increased clear discharge may occur for a few days. Patients should receive individualized instructions about activity and vaginal intercourse; routine prolonged bed rest or activity restriction is not recommended solely to prevent preterm birth. They should seek urgent assessment for regular contractions, heavy bleeding, fluid leakage, fever, foul-smelling discharge, or severe pain.

Follow-up is based on the indication and clinical course. Routine serial cervical-length surveillance after a cerclage is not required for every patient. A transvaginal cerclage is generally removed at about 37 weeks, or earlier if labor, PPROM, significant bleeding, or infection develops. A transabdominal cerclage is usually left in place and delivery occurs by cesarean birth, commonly between 37\,0/7 and 39\,0/7 weeks when there is no earlier indication for delivery.

\section*{Outcomes \& Prognosis}
Benefit depends strongly on the indication, gestational age, cervical findings, prior obstetric history, and competing causes of preterm birth. Cerclage can reduce recurrent preterm birth in selected high-risk patients, but it cannot prevent all preterm births and should not be described with a single universal success rate. Future pregnancies require an individualized plan based on the prior indication and outcome.

\section*{Clinical Decision-Making and Counseling}
The indication should be separated from the ultrasound finding itself. A short cervix is a measurement, not a diagnosis of cervical insufficiency. In an asymptomatic singleton pregnancy without a prior spontaneous preterm birth, vaginal progesterone and surveillance are commonly considered, while cerclage is not routinely placed solely because the cervix is short. A patient with a relevant prior spontaneous preterm birth and a newly short cervix may meet criteria for an ultrasound-indicated stitch, depending on gestational age and the degree of shortening. A history-indicated stitch is based on the pattern of the previous loss or early birth, especially painless dilation, rather than on a single prior preterm birth from another cause such as preeclampsia or placental abruption.

Before an exam-indicated stitch, the clinician should establish that the membranes are intact and that the presentation is compatible with continuing the pregnancy. Fever, uterine tenderness, foul discharge, regular contractions, significant bleeding, or fetal compromise should prompt evaluation for infection, labor, placental pathology, or another reason not to proceed. Amniocentesis for infection testing is not required in every case, but may be considered when infection is suspected and the result would change management. The possible benefit of rescue cerclage must be balanced against ruptured membranes, infection, hemorrhage, cervical trauma, and the possibility that the procedure will not prolong the pregnancy.

The stitch is a support, not a mechanical guarantee that the cervix will remain closed. Patients should understand that preterm birth can result from uterine contractions, membrane rupture, placental disease, fetal indications, infection, or causes unrelated to cervical dilation. The procedure does not replace treatment of infection, progesterone when indicated, smoking cessation, or management of other obstetric risks. Routine bed rest is not recommended solely to prevent preterm birth; activity recommendations should be individualized.

Transabdominal cerclage deserves specific counseling because it generally cannot be removed vaginally. It may be placed before pregnancy or in early pregnancy, through open or minimally invasive surgery, and usually necessitates cesarean delivery. It can be considered after a prior transvaginal cerclage failed at an early gestational age or when the cervix is too short or surgically altered for a safe vaginal stitch. Counseling should include operative risks, the need for cesarean birth, future pregnancy planning, and the possibility of leaving the stitch in place for subsequent pregnancies.

After placement, routine digital examinations are not needed unless symptoms occur. The patient should know the expected degree of spotting and cramping, the location of the stitch, and how to access urgent obstetric care. If membranes rupture, contractions begin, or infection is suspected, the stitch is assessed promptly; leaving it in place during active labor can cause cervical laceration, while removing it too early can forfeit any remaining benefit. At removal, the clinician documents the cervical appearance, bleeding, membrane status, and fetal or maternal indication. The subsequent pregnancy plan should be reviewed postpartum, when operative records and the gestational age at delivery can guide future surveillance or cerclage.

\section*{Postpartum Review}
The indication and outcome should be reviewed after delivery. The record should state whether the stitch was history-, ultrasound-, exam-, or transabdominal-indicated, the gestational age at placement and removal, membrane status, complications, and the gestational age and reason for delivery. This information is more useful for planning a future pregnancy than labeling the cerclage a simple success or failure. A subsequent pregnancy may require early maternal--fetal medicine review, cervical-length surveillance, progesterone, or repeat cerclage according to the prior course.

Patients should be counseled that a cerclage does not imply that the cervix is permanently weak, and a future pregnancy is not automatically managed identically. Conversely, a prior uncomplicated term delivery after a stitch does not eliminate all future risk. Interpregnancy health, smoking cessation, chronic-disease control, and review of the prior records are part of prevention planning.

\section*{Documentation and Safety}
The operative note records the indication, gestational age, cervical appearance, stitch material and position, complications, fetal assessment, and removal plan. Clear documentation is important because future clinicians may need to decide whether a stitch is present before labor, membrane rupture, cervical examination, or delivery. The patient should be given the procedure details and emergency contact instructions.

\section*{Practical Follow-Up}
Cerclage does not remove the need for routine prenatal care, fetal surveillance when indicated, infection prevention, or evaluation of contractions and bleeding. The patient should know the planned removal date and whether cesarean birth is required. Any change in symptoms should be assessed by the obstetric team rather than managed by self-imposed bed rest.

\section*{Safety Summary}
Cerclage is offered only when the likely benefit outweighs the risks for that pregnancy. It does not replace assessment for labor, membrane rupture, infection, placental disease, or other causes of preterm birth, and its removal plan must be clear before delivery.

\section*{Final Safety Note}
Cerclage is not a substitute for assessing contractions, membrane rupture, infection, bleeding, placental disease, or fetal indications. The stitch-removal and delivery plan should be documented and communicated to the patient.

\section*{Completion Checklist}
Document the indication, gestational age, stitch type and location, removal plan, membrane status, and warning symptoms. The patient should know when to seek urgent care and which delivery route may be required.

\section*{Handoff}
The obstetric handoff should identify the stitch and removal plan before any presentation in labor, rupture of membranes, bleeding, or planned delivery.

\section*{References}
\begin{enumerate}
  \item American College of Obstetricians and Gynecologists. Cerclage for the Management of Cervical Insufficiency. Practice Bulletin No. 142. Reaffirmed 2024.
  \item American College of Obstetricians and Gynecologists. Cervical Cerclage. FAQ526. Updated 2024; reviewed 2025.
  \item Society for Maternal--Fetal Medicine. Consult Series No. 70: Management of Short Cervix in Individuals Without a History of Spontaneous Preterm Birth. 2024.
  \item Society for Maternal--Fetal Medicine. Consult Series No. 65: Transabdominal Cerclage. 2023; reaffirmed 2025.
\end{enumerate}

\clearpage
